\section{Experiments: Natural Language Processing} \label{sec:expt-nlp}

In this section our goal is to showcase benefits of modeling longer input sequence for NLP tasks, for which we select three representative tasks.
We begin with basic masked language modeling (MLM; \citealt{devlin2018bert}) to check if better contextual representations can be learnt by utilizing longer contiguous sequences.
Next, we consider QA with supporting evidence, for which capability to handle longer sequence would allow us to retrieve more evidence using crude systems like TF-IDF/BM25.
Finally, we tackle long document classification where discriminating information may not be located in first 512 tokens.
Below we summarize the results for \bigb using sequence length 4096\footnote{code available at \url{http://goo.gle/bigbird-transformer}}, while we defer all other setup details including computational resources, batch size, step size, to~\Cref{sec:apndx-expt-nlp}. 


\paragraph{Pretraining and MLM}
We follow \citep{devlin2018bert, liu2019roberta} to create base and large versions of \bigb and pretrain it using MLM objective. 
This task involves predicting a random subset of tokens which have been masked out. 
We use four standard data-sets for pretraining (listed in \Cref{sec:app-expt-nlp:mlm}, \Cref{tab:mlm_data}), warm-starting from the public RoBERTa checkpoint\footnote{\url{https://github.com/pytorch/fairseq/tree/master/examples/roberta}}. 
We compare performance in predicting the masked out tokens in terms of bits per character, following~\citep{beltagy2020longformer}. 
As seen in~\Cref{sec:app-expt-nlp:mlm}, ~\Cref{tab:mlm_bpc},  both \bigb and Longformer perform better than limited length RoBERTa, with \bigb-\textsc{etc} performing the best.
We note that we trained our models on a 
reasonable $16GB$ memory/chip with batch size of 32-64.
Our memory efficiency is due to efficient blocking and sparsity structure of 
the sparse attention mechanism described in~\Cref{sec:arch}. 
\begin{table}
    \centering
    \small
    \begin{tabular}{@{}l c c c ccc c cc c c@{}}
    \toprule
    \multirow{2}{*}{Model} & & 
    \multicolumn{3}{c}{HotpotQA} & &
    \multicolumn{2}{c}{NaturalQ} & &
    TriviaQA & & 
    WikiHop\\
    \cmidrule{3-5} \cmidrule{7-8} \cmidrule{10-10}  \cmidrule{12-12} 
               && Ans & Sup & Joint  &&  LA & SA && Full && MCQ  \\
    \midrule
    RoBERTa     && 73.5 & 83.4 & 63.5 &&  -   &  -  && 74.3 && 72.4  \\
    Longformer  && 74.3 & 84.4 & 64.4 &&  -   &  -  && 75.2 && 75.0  \\
    \bigb-\textsc{itc}  && \textbf{75.7} & 86.8 & 67.7 && 70.8 & 53.3 && \textbf{79.5} && \textbf{75.9} \\
    \bigb-\textsc{etc}  && 75.5 & \textbf{87.1} & \textbf{67.8} && \textbf{73.9} & \textbf{54.9} && 78.7 && \textbf{75.9} \\
    \bottomrule
    \end{tabular}
    \vspace{2mm}
    \caption{QA Dev results using Base size models. We report accuracy for WikiHop and F1 for HotpotQA, Natural Questions, and TriviaQA.}
    \label{tab:QADev}
\end{table}

\begin{table}
    \centering
    \small
    \begin{tabular}{@{}l c c ccc c cc c cc@{}}
    \toprule
    \multirow{2}{*}{Model} & 
    \multicolumn{3}{c}{HotpotQA} & &
    \multicolumn{2}{c}{NaturalQ} & &
    \multicolumn{2}{c}{TriviaQA} & &
    WikiHop\\
    \cmidrule{2-4} \cmidrule{6-7} \cmidrule{9-10} \cmidrule{12-12} 
    & Ans & Sup & Joint  &&  LA & SA && Full & Verified && MCQ  \\
    \midrule
    HGN \citep{fang2019hierarchical}         & \textbf{82.2} & 88.5 & \textbf{74.2} &&  -   &  -   &&  -   &  -   &&  -   \\
    GSAN                                     & 81.6 & 88.7 & 73.9 &&  -   &  -   &&  -   &  -   &&  -   \\
    ReflectionNet \citep{gong2020reflection} &  -   &  -   &  -   && 77.1 & \textbf{64.1} &&  -   &  -   &&  -   \\
    RikiNet-v2 \citep{liu2020rikinet}           &  -   &  -   &  -   && 76.1 & 61.3 &&  -   &  -   &&  -   \\
    Fusion-in-Decoder \citep{izacard2020fid} &  -   &  -   &  -   &&  -   &  -   && 84.4 & 90.3 &&  -   \\
    SpanBERT \citep{joshi2020spanbert}       &  -   &  -   &  -   &&  -   &  -   && 79.1 & 86.6 &&  -   \\ 
    MRC-GCN \citep{tang2020multi}            &  -   &  -   &  -   &&  -   &  -   &&  -   &  -   && 78.3 \\
    MultiHop \citep{chen2019multi}           &  -   &  -   &  -   &&  -   &  -   &&  -   &  -   && 76.5 \\
    Longformer \citep{beltagy2020longformer} & 81.2 & 88.3 & 73.2 &&  -   &  -   && 77.3 & 85.3 && 81.9 \\
    \midrule
    \bigb-\textsc{etc} & 81.2 & \textbf{89.1} & 73.6 && \textbf{77.8} & 57.9 && \textbf{84.5} & \textbf{92.4} && \textbf{82.3} \\
    \bottomrule
    \end{tabular}
    \vspace{2mm}
    \caption{Fine-tuning results on \textbf{Test} set for QA tasks. The Test results (F1 for HotpotQA, Natural Questions, TriviaQA, and Accuracy for WikiHop) have been picked from their respective leaderboard. For each task the top-3 leaders were picked not including \bigb-etc. \textbf{For Natural Questions Long Answer (LA), TriviaQA, and WikiHop, \bigb-ETC is the new state-of-the-art}. On HotpotQA we are third in the leaderboard by F1 and second by Exact Match (EM).}
    \label{tab:QATest}
\end{table}


\paragraph{Question Answering (QA)} 
We considered following four challenging datasets:
\vspace{-1mm}
\begin{enumerate}[leftmargin=6mm, itemsep=2mm, partopsep=0pt,parsep=0pt]
\item  Natural Questions~\citep{kwiatkowski2019natural}: 
For the given question, find a short span of answer (SA) from the given evidences as well highlight the paragraph from the given evidences 
containing information about the correct answer (LA).
\item HotpotQA-distractor~\citep{yang2018hotpotqa}: Similar to natural questions, it requires finding the answer (Ans) as well as the supporting facts (Sup) over different documents needed for multi-hop reasoning from the given evidences. 
\item TriviaQA-wiki~\citep{JoshiTriviaQA2017}: We need to provide an answer for the given question using provided Wikipedia evidence, however, the answer might not be present in the given evidence. On a smaller \emph{verified} subset of question, the given evidence is guaranteed to contain the answer. Nevertheless, we model the answer as span selection problem in this case as well.
\item  WikiHop~\citep{welbl2018constructing}: Chose correct option from multiple-choice questions (MCQ), by aggregating information spread across multiple documents given in the evidences.
\end{enumerate}
As these tasks are very competitive, multiple highly engineered systems have been designed specific each dataset confirming to respective output formats.
For a fair comparison, we had to use some additional regularization for training \bigb, details of which are provided in~\Cref{sec:app-expt-nlp:qa} along with exact architecture description.
We experiment using the base sized model and select the best configuration on the development set for each dataset (as reported in~\Cref{tab:QADev}).
We can see that \bigb-\textsc{etc}, with expanded global tokens consistently outperforms all other models.
Thus, we chose this configuration to train a large sized model to be used for evaluation on the hidden test set.

In~\Cref{tab:QATest}, we compare \bigb-\textsc{etc} model to top-3 entries from the leaderboard excluding \bigb. 
One can clearly see the importance of using longer context as both Longformer and \bigb outperform models with smaller contexts.
Also, it is worth noting that \bigb submission is a single model, whereas the other top-3 entries for Natural Questions are  ensembles, which might explain the slightly lower accuracy in exact answer phrase selection.


\paragraph{Classification} 
We experiment on datasets of different lengths and contents, specifically various document classification and GLUE tasks.
Following BERT, we used one layer with cross entropy loss on top of the first [CLS] token.
We see that gains of using \bigb are more significant when we have longer documents and fewer training examples.
For instance, using base sized model,
\bigb improves state-of-the-art for Arxiv dataset by about $\bm{5\%}$ \textbf{points}.
On Patents dataset, there is improvement over using simple BERT/RoBERTa, but given the large size of training data the improvement over SoTA (which is not BERT based) is not significant.
Note that this performance gain is not seen for much
smaller IMDb dataset.
Along with experimental setup detail, we present detailed results in~\Cref{sec:app-expt-nlp:cls} which show competitive performance.


\begin{table}[b]
    \centering
    \small
    % \hspace*{-7mm}
    \begin{tabular}{@{}p{1mm}l @{}p{3mm}@{} rrr @{}p{5mm}@{} rrr @{}p{5mm}@{} rrr@{}}
    \toprule
     \multicolumn{2}{l}{\multirow[b]{2}{*}{\hspace{-2mm}\normalsize{Model}}} & &
     \multicolumn{3}{c}{Arxiv} & & \multicolumn{3}{c}{PubMed} & & \multicolumn{3}{c}{BigPatent}\\
    \cmidrule{4-6} \cmidrule{8-10} \cmidrule{12-14}
    & & & R-1 & R-2 & R-L & & R-1 & R-2 & R-L & & R-1 & R-2 & R-L \\
    \midrule
    \multirow{10}{*}{\rotatebox[origin=c]{90}{Prior Art}}
    & SumBasic~\citep{nenkova2005impact}          & & 29.47 &  6.95 & 26.30 & & 37.15 & 11.36 & 33.43 & & 27.44 & 7.08 & 23.66\\
    & LexRank~\citep{erkan2004lexrank}          & & 33.85 & 10.73 & 28.99 & & 39.19 & 13.89 & 34.59 & & 35.57 & 10.47 & 29.03 \\
    & LSA~\citep{wiseman2017challenges}               & & 29.91 &  7.42 & 25.67 & & 33.89 &  9.93 & 29.70 & & - & - & - \\
    & Attn-Seq2Seq~\citep{sutskever2014sequence}    & & 29.30 &  6.00 & 25.56 & & 31.55 &  8.52 & 27.38 & & 28.74 & 7.87 & 24.66 \\
    & Pntr-Gen-Seq2Seq~\citep{see2017get} & & 32.06 &  9.04 & 25.16 & & 35.86 & 10.22 & 29.69 & &  33.14 & 11.63 & 28.55 \\
    & Long-Doc-Seq2Seq~\citep{cohan2018discourse} & & 35.80 & 11.05 & 31.80 & & 38.93 & 15.37 & 35.21 & & - & - & - \\
    & Sent-CLF~\citep{subramanian2019extractive}  & & 34.01 &  8.71 & 30.41 & & 45.01 & 19.91 & 41.16 & & 36.20 & 10.99 & 31.83 \\
    & Sent-PTR~\citep{subramanian2019extractive}  & & 42.32 & 15.63 & 38.06 & & 43.30 & 17.92 & 39.47 & & 34.21 & 10.78 & 30.07 \\
    & Extr-Abst-TLM~\citep{subramanian2019extractive} & & 41.62 & 14.69 & 38.03 & & 42.13 & 16.27 & 39.21 & & 38.65 & 12.31 & 34.09 \\
    & Dancer~\citep{gidiotis2020divide}  & & 42.70 & 16.54 & 38.44 & & 44.09 & 17.69 & 40.27 & & - & - & - \\
    \midrule
    \multirow{4}{*}{\rotatebox[origin=c]{90}{Base}}
    & Transformer & & 28.52 &  6.70 & 25.58 & & 31.71 &  8.32 & 29.42 & & 39.66 & 20.94 & 31.20 \\
    & \; + RoBERTa~\citep{rothe2019leveraging} & & 31.98 &  8.13 & 29.53  & & 35.77 & 13.85 & 33.32 & & 41.11 & 22.10 & 32.58 \\
    & \; + Pegasus~\citep{zhang2019pegasus}  & & 34.81 & 10.16 & 30.14 & & 39.98 & 15.15 & 35.89 & & 43.55 & 20.43 & 31.80 \\
    & \bigb-RoBERTa & & \underline{41.22} & \underline{16.43} & \underline{36.96} & & \underline{43.70} & \underline{19.32} & \underline{39.99} & & \underline{55.69} & \underline{37.27} & \underline{45.56} \\
    \midrule
    \multirow{3}{*}{\rotatebox[origin=c]{90}{Large}}
    & Pegasus (Reported)~\citep{zhang2019pegasus} & & 44.21 & 16.95 & 38.83 & & 45.97 & 20.15 & 41.34 & & 52.29 & 33.08 & 41.75 \\
    & Pegasus (Re-eval)                          & & 43.85 & 16.83 & 39.17 & & 44.53 & 19.30 & 40.70 & & 52.25 & 33.04 & 41.80 \\
    & \bigb-Pegasus & & \textbf{46.63} & \textbf{19.02} & \textbf{41.77}   & & \textbf{46.32} & \textbf{20.65} & \textbf{42.33}   & & \textbf{60.64} & \textbf{42.46} & \textbf{50.01}  \\
    \bottomrule
    \end{tabular}
    \vspace{2mm}
    \caption{Summarization ROUGE score for long documents.}
    \label{tab:long_sum_res}
\end{table}


\subsection{Encoder-Decoder Tasks}
\label{sec:seq2seq}
For an encoder-decoder setup, one can easily see that both suffer from quadratic complexity due to the full self attention.
We focus on introducing the sparse attention mechanism of \bigb only at the encoder side.
This is because, in practical generative applications, the length of output sequence is typically small as compared to the input. 
For example for text summarization, we see in realistic scenarios (c.f.~\Cref{sec:appn_summarization}~\Cref{tab:long_sum_data}) that the median output sequence length is $\sim 200$ where as the input sequence's median length is $>3000$.
For such applications, it is more efficient to use sparse attention mechanism for the encoder and full self-attention for the decoder.

\paragraph{Summarization}
Document summarization is a task of creating a short and accurate summary of a text document. 
We used three long document datasets for testing our model details of which are mention in~\Cref{tab:long_sum_data}.
In this paper we focus on abstractive summarization of long documents where using a longer contextual encoder should improve performance.
The reasons are two fold:
First, the salient content can be evenly distributed in the long document, not just in first 512 tokens, and this is by design in the BigPatents dataset~\citep{sharma2019bigpatent}.
Second, longer documents exhibit a richer discourse structure and summaries are considerably more abstractive, thereby observing more context helps.
As has been pointed out recently ~\citep{rothe2019leveraging,zhang2019pegasus}, pretraining helps in generative tasks, we warm start from our general purpose MLM pretraining on base-sized models as well as utilizing state-of-the-art summarization specific pretraining from Pegasus~\citep{zhang2019pegasus} on large-sized models.
The results of training \bigb sparse encoder along with full decoder on these long document datasets are presented in~\Cref{tab:long_sum_res}.
We can clearly see modeling longer context brings significant improvement. 
Along with hyperparameters, we also present results on shorter but more widespread datasets in~\Cref{sec:appn_summarization}, which show that using sparse attention does not hamper performance either.

